\title{The Era of 1-bit LLMs: All Large Language Models are in 1.58 Bits}

\begin{abstract}
Advancements in research exemplified by BitNet~\cite{bitnet} have initiated a new epoch characterized by 1-bit Large Language Models (LLMs). In this paper, we present a novel 1-bit LLM, designated as \textbf{\text{BitNet b1.58}}, where each parameter (or weight) assumes one of the ternary values from the set \{-1, 0, 1\}. Remarkably, this model attains parity with its full-precision (FP16 or BF16) Transformer LLM counterpart, maintaining comparable perplexity and downstream performance for an equivalent model size and number of training tokens, but with notable improvements in latency, memory utilization, throughput, and energy efficiency. Beyond its cost-effectiveness, the 1.58-bit LLM establishes a new scaling paradigm and methodological framework for training subsequent high-performance and resource-efficient LLMs. Additionally, this development paves the way for innovative computing architectures and offers the prospect of developing hardware tailored specifically to optimize 1-bit LLM operations.
\end{abstract}

\section{The Era of 1-bit LLMs}

Artificial intelligence research has witnessed a significant expansion in the magnitude and capability of Large Language Models (LLMs) in recent years. These systems have achieved state-of-the-art results across numerous natural language processing benchmarks, yet their escalating scale presents formidable obstacles to practical deployment and increases apprehensions related to environmental sustainability and financial viability, primarily driven by substantial energy demands.

A promising direction to mitigate these issues is post-training quantization, wherein models are compressed to utilize reduced-bit representations during inference~\cite{smoothquant, gptq, quip, quip_sharp}. By decreasing the bitwidth of model weights and activations, this strategy offers pronounced benefits in memory consumption and computational burden for LLMs. Recent progressions include lowering the precision to 4-bit formats~\cite{gptq, awq}. Nevertheless, despite its prevalence in industrial applications, post-training quantization remains a sub-optimal solution.

Emerging research into 1-bit model architectures—exemplified by BitNet~\cite{bitnet}—provides a compelling path toward lowering large language model (LLM) costs while preserving accuracy. Traditional LLMs operate on 16-bit floating point formats (FP16 or BF16), with the computational burden dominated by matrix multiplication, making floating-point additions and multiplications the primary contributors to resource consumption. In contrast, BitNet’s matrix multiplications employ only integer additions, drastically decreasing the associated energy usage. Because energy consumption often constrains compute performance on modern hardware, these savings can directly enable higher throughput.

Beyond computational considerations, the transfer of model weights from DRAM to the accelerator’s on-chip memory (e.g., SRAM) incurs significant inference costs. While increasing SRAM capacity is one approach to boosting throughput, it comes with substantial expense relative to DRAM. By reducing the precision to 1-bit, LLMs substantially shrink their memory footprint in both size and required bandwidth, lowering the expense and latency of weight transfers and thereby facilitating faster, more cost-effective inference.

This paper introduces \textbf{\text{BitNet b1.58}}, an enhanced 1-bit LLM model where every weight is ternary, drawn from \{-1, 0, 1\}. Extending the standard BitNet, the addition of the 0 value yields 1.58 bits per parameter. \text{BitNet b1.58} upholds the original 1-bit BitNet’s efficiency, including its computation strategy that essentially eliminates multiplications in matrix operations and is amenable to aggressive optimization. Its power and memory demands remain low, leading to improved throughput and lower latency versus FP16-based models. Importantly, \text{BitNet b1.58} offers two distinct advantages: (1) it features enhanced modeling power through explicit feature filtering, enabled by permitting 0-valued weights—a modification that can notably improve the effectiveness of 1-bit LLMs; and (2) empirical results indicate that, with the same settings (e.g., model parameters, token count), \text{BitNet b1.58} can achieve parity with full-precision (FP16) models in terms of perplexity and downstream task performance, starting at the 3B parameter scale.

\section{BitNet b1.58}

\text{BitNet b1.58} is built upon the BitNet framework, a type of Transformer that substitutes \emph{nn.Linear} layers with \emph{BitLinear}. This model is trained from the beginning, utilizing weights quantized to 1.58 bits and activations quantized to 8 bits. In contrast to the original BitNet, this iteration incorporates several adjustments, which we outline below.

\paragraph{Quantization Function.} To enforce that weight values are limited to -1, 0, or +1, we employ an \emph{absmean} quantization scheme. This approach normalizes the weight matrix by dividing by its mean absolute value, then rounds each entry to the closest integer within \{-1, 0, +1\}:

\begin{equation}
    \widetilde{W} = \mathrm{RoundClip}(\frac{W}{\gamma + \epsilon}, -1, 1),
\end{equation}

\begin{equation}
    \mathrm{RoundClip}(x, a, b) = \max(a, \min(b, \mathrm{round}(x))),
\end{equation}

\begin{equation}
    \gamma = \frac{1}{nm}\sum_{ij} |W_{ij}|.
\end{equation}

The activation quantization mechanism is consistent with that in BitNet, except for a key difference: we omit scaling activations into the interval $[0, Q_b]$ prior to applying non-linearities. Instead, each activation is rescaled to fit within $[-Q_b, Q_b]$ per token, thereby eliminating the need for zero-point quantization. This change simplifies both implementation and system-level tuning, while we observe minimal impact on empirical performance.

\paragraph{LLaMA-alike Components.}

The LLaMA~\cite{llama, llama2} architecture has become the standard foundation for open-source large language models (LLMs). To support the open-source ecosystem, we incorporate LLaMA-style components into \text{BitNet b1.58}. In particular, our model leverages RMSNorm~\cite{rmsnorm}, SwiGLU~\cite{swiglu}, and rotary positional embedding~\cite{rope}, and removes bias terms throughout. These choices allow \text{BitNet b1.58} to be seamlessly compatible with major open-source platforms (such as Huggingface, vLLM~\cite{vllm}, and llama.cpp\footnote{\url{https://github.com/ggerganov/llama.cpp}}) with minimal modification.

\section{Results}

We conducted a comparison between \text{BitNet b1.58} and our reproduced FP16 \text{LLaMA LLM} across multiple model scales. Both models underwent pre-training on the RedPajama dataset~\cite{redpajama} for 100 billion tokens to maintain experimental parity. Their zero-shot capabilities were assessed on a spectrum of language understanding tasks, encompassing ARC-Easy~\cite{arc}, ARC-Challenge~\cite{arc}, Hellaswag~\cite{hellaswag}, Winogrande~\cite{winoGrande}, PIQA~\cite{piqa}, OpenbookQA~\cite{openbookqa}, and BoolQ~\cite{boolq}. Additionally, we recorded validation perplexity scores on WikiText2~\cite{wikitext2} and C4~\cite{c4} corpora.

To assess efficiency, both GPU memory usage and latency were measured for \text{LLaMA LLM} and \text{BitNet b1.58} during inference. These benchmarks were obtained via the FasterTransformer\footnote{\url{https://github.com/NVIDIA/FasterTransformer}} framework, which is highly tuned for reducing LLM inference latency on GPUs. The Ladder~\cite{ladder} 2-bit kernel was utilized specifically for \text{BitNet b1.58}. Time per output token was used as the primary inference performance metric.

In Table~\ref{tab:ppl}, we present perplexity alongside resource consumption metrics for both \text{BitNet b1.58} and \text{LLaMA LLM}. Notably, \text{BitNet b1.58} reaches the perplexity of its full-precision counterpart starting at the 3B scale, while achieving 2.71$\times$ higher speed and requiring 3.55$\times$ less GPU memory. At the 3.9B parameter scale, \text{BitNet b1.58} achieves a 2.4$\times$ acceleration, reduces memory demands by 3.32$\times$, and surpasses the accuracy of \text{LLaMA LLM} 3B.

Table~\ref{tab:zero-shot} details the zero-shot accuracy outcomes across downstream tasks, following the evaluation protocol implemented in \emph{lm-evaluation-harness}\footnote{\url{https://github.com/EleutherAI/lm-evaluation-harness}}. The findings reveal that as model sizes grow, the discrepancy in accuracy between \text{BitNet b1.58} and \text{LLaMA LLM} decreases. Crucially, performance parity with the full-precision baseline emerges at the 3B scale. Corroborating the perplexity trends, the task results indicate that \text{BitNet b1.58} at 3.9B consistently surpasses the \text{LLaMA LLM} 3B model in both efficiency and performance, reinforcing the conclusion that \text{BitNet b1.58} offers a Pareto improvement compared to leading LLM architectures.

\paragraph{Memory and Latency}

The model was further expanded to sizes of 7B, 13B, and 70B, and we assessed the associated costs. As depicted in Figure~\ref{fig:latency-memory-energy}, both latency and memory usage exhibit consistent patterns: the acceleration effect grows with increasing model scale. Notably, \text{BitNet b1.58} 70B achieves a 4.1-fold improvement in speed relative to the \text{LLaMA LLM} reference. This performance gain is attributed to the scaling of time expenditure in \emph{nn.Linear} operations as model size increases. Memory utilization also follows this pattern; since the embeddings remain at full precision, their proportional memory requirement decreases in larger architectures. All latency and memory measurements utilize a 2-bit kernel, suggesting further reductions are possible through additional optimization.

\paragraph{Energy}

We estimated the energy consumption due to arithmetic operations for both \text{BitNet b1.58} and \text{LLaMA LLM}. Our primary focus was the matrix multiplication operations, as they constitute the bulk of computation in LLMs. Figure~\ref{fig:energy} provides a breakdown of the energy consumption. For \text{BitNet b1.58}, INT8 additions dominate, whereas \text{LLaMA LLM} requires both FP16 multiplications and additions. Using the energy analysis framework described in~\cite{energycost, pokebnn}, our calculations indicate that \text{BitNet b1.58} lowers the arithmetic operation energy needed for matrix multiplication by a factor of 71.4 when deployed on 7nm hardware. Additionally, we reported comprehensive energy usage for tasks involving 512 tokens. The findings demonstrate that as model size increases, the energy efficiency advantage of \text{BitNet b1.58} over the FP16 \text{LLaMA LLM} becomes more pronounced. This arises because the proportion of time and energy devoted to \emph{nn.Linear} operations grows with model scale, while other elements contribute less substantially in larger models.

\paragraph{Throughput}

We evaluate and contrast the throughput of \text{BitNet b1.58} and \text{LLaMA LLM} with 70B parameters, utilizing two 80GB A100 GPUs and employing pipeline parallelism~\cite{gpipe} to enable the execution of \text{LLaMA LLM} 70B on this hardware. The batch size was increased incrementally up to the maximum permitted by GPU memory, with a fixed sequence length of 512. As shown in Table~\ref{tab:throughput}, \text{BitNet b1.58} 70B is capable of handling a batch size up to 11 times greater than that of \text{LLaMA LLM}, and this advantage yields an 8.9-fold improvement in throughput.

\textbf{\text{BitNet b1.58} introduces a novel scaling relationship between model performance and inference efficiency.} For context, Figures~\ref{fig:latency-memory-energy} and \ref{fig:energy} establish the following correspondences in terms of 1.58-bit versus 16-bit models:

\begin{itemize}
    \item The 13B BitNet b1.58 configuration outperforms the 3B FP16 LLM regarding latency, memory demands, and energy requirements.
    \item The 30B BitNet b1.58 model demonstrates higher efficiency across latency, memory usage, and energy metrics when compared to the 7B FP16 LLM.
    \item The 70B BitNet b1.58 variant surpasses the 13B FP16 LLM in terms of latency, memory efficiency, and energy savings.
\end{itemize}

\paragraph{Training with 2T Tokens}

Training token count is pivotal for scaling LLMs. To examine the token-level scalability of \text{BitNet b1.58}, we trained a \text{BitNet b1.58} model on 2T tokens using the data curation pipeline outlined for StableLM-3B~\cite{StableLM-3B-4E1T}, which represents a state-of-the-art 3B open-source model. Both models were assessed using a composite benchmark encompassing Winogrande~\cite{winoGrande}, PIQA~\cite{piqa}, SciQ~\cite{sciq}, LAMBADA~\cite{lambada}, and ARC-easy~\cite{arc}. The zero-shot results are reported in Table~\ref{tab:2t}. For benchmarks with both accuracy and normalized accuracy, their mean was used for comparison. StableLM 3b results at 2T tokens were sourced from its official report. Our analysis reveals that \text{BitNet b1.58} achieves higher performance across all evaluated tasks, supporting the conclusion that 1.58-bit LLMs exhibit strong generalization.

\section{Discussion and Future Work}

\textbf{1-bit Mixture-of-Experts (MoE) LLMs}

Mixture-of-Experts (MoE) has emerged as a resource-efficient strategy for scaling LLMs, significantly reducing the floating-point operations required for inference. However, MoE's practical use is frequently hampered by substantial memory requirements and intensive communication between computational nodes. Employing 1.58-bit LLMs offers solutions to these bottlenecks. The smaller memory demand permits larger models to fit onto fewer devices, and also minimizes the network traffic associated with moving activations. Ideally, placing an entire MoE on a single device would eliminate communication overhead completely.

\textbf{Native Support of Long Sequence in LLMs}

Long-context modeling has become essential for contemporary LLM applications. A persistent obstacle is the extensive memory usage required by KV caches during long sequence inference. \text{BitNet b1.58} makes significant progress in this regard by halving the activation precision from 16 bits to 8 bits, effectively enabling double the context length within the same resource envelope. This activation footprint could be further reduced to 4 bits—or potentially even lower—if extended to 1.58-bit LLMs, which presents an interesting avenue for future research.

\textbf{LLMs on Edge and Mobile}

Deploying 1.58-bit LLMs presents promising opportunities for edge and mobile scenarios, where limitations in computational power and memory typically constrain model size and performance. The low memory and reduced energy requirements of 1.58-bit LLMs make it feasible to implement LLMs in these settings, thereby expanding the range of achievable applications on edge and mobile platforms. Additionally, these bit-width LLMs are well-aligned with the capabilities of CPUs, the predominant processors for such devices, allowing \text{BitNet b1.58} to operate efficiently and unlocking improved device functionality.

\textbf{New Hardware for 1-bit LLMs}

Cutting-edge efforts, such as those demonstrated by Groq\footnote{\url{https://groq.com/}}, illustrate the feasibility and benefits of developing specialized hardware—like LPUs—tailored to the demands of LLM workloads. We foresee further innovation in hardware design explicitly targeted at the requirements of 1-bit LLMs, leveraging the novel computational efficiencies realized through BitNet~\cite{bitnet} and advocating for future system-level advancements in this space.

\end{document}